% Options for packages loaded elsewhere
\PassOptionsToPackage{unicode}{hyperref}
\PassOptionsToPackage{hyphens}{url}
%
\documentclass[
]{scrbook}

\usepackage{amsmath,amssymb}
\usepackage{iftex}
\ifPDFTeX
  \usepackage[T1]{fontenc}
  \usepackage[utf8]{inputenc}
  \usepackage{textcomp} % provide euro and other symbols
\else % if luatex or xetex
  \usepackage{unicode-math}
  \defaultfontfeatures{Scale=MatchLowercase}
  \defaultfontfeatures[\rmfamily]{Ligatures=TeX,Scale=1}
\fi
\usepackage{lmodern}
\ifPDFTeX\else  
    % xetex/luatex font selection
\fi
% Use upquote if available, for straight quotes in verbatim environments
\IfFileExists{upquote.sty}{\usepackage{upquote}}{}
\IfFileExists{microtype.sty}{% use microtype if available
  \usepackage[]{microtype}
  \UseMicrotypeSet[protrusion]{basicmath} % disable protrusion for tt fonts
}{}
\makeatletter
\@ifundefined{KOMAClassName}{% if non-KOMA class
  \IfFileExists{parskip.sty}{%
    \usepackage{parskip}
  }{% else
    \setlength{\parindent}{0pt}
    \setlength{\parskip}{6pt plus 2pt minus 1pt}}
}{% if KOMA class
  \KOMAoptions{parskip=half}}
\makeatother
\usepackage{xcolor}
\setlength{\emergencystretch}{3em} % prevent overfull lines
\setcounter{secnumdepth}{1}
% Make \paragraph and \subparagraph free-standing
\makeatletter
\ifx\paragraph\undefined\else
  \let\oldparagraph\paragraph
  \renewcommand{\paragraph}{
    \@ifstar
      \xxxParagraphStar
      \xxxParagraphNoStar
  }
  \newcommand{\xxxParagraphStar}[1]{\oldparagraph*{#1}\mbox{}}
  \newcommand{\xxxParagraphNoStar}[1]{\oldparagraph{#1}\mbox{}}
\fi
\ifx\subparagraph\undefined\else
  \let\oldsubparagraph\subparagraph
  \renewcommand{\subparagraph}{
    \@ifstar
      \xxxSubParagraphStar
      \xxxSubParagraphNoStar
  }
  \newcommand{\xxxSubParagraphStar}[1]{\oldsubparagraph*{#1}\mbox{}}
  \newcommand{\xxxSubParagraphNoStar}[1]{\oldsubparagraph{#1}\mbox{}}
\fi
\makeatother


\providecommand{\tightlist}{%
  \setlength{\itemsep}{0pt}\setlength{\parskip}{0pt}}\usepackage{longtable,booktabs,array}
\usepackage{calc} % for calculating minipage widths
% Correct order of tables after \paragraph or \subparagraph
\usepackage{etoolbox}
\makeatletter
\patchcmd\longtable{\par}{\if@noskipsec\mbox{}\fi\par}{}{}
\makeatother
% Allow footnotes in longtable head/foot
\IfFileExists{footnotehyper.sty}{\usepackage{footnotehyper}}{\usepackage{footnote}}
\makesavenoteenv{longtable}
\usepackage{graphicx}
\makeatletter
\newsavebox\pandoc@box
\newcommand*\pandocbounded[1]{% scales image to fit in text height/width
  \sbox\pandoc@box{#1}%
  \Gscale@div\@tempa{\textheight}{\dimexpr\ht\pandoc@box+\dp\pandoc@box\relax}%
  \Gscale@div\@tempb{\linewidth}{\wd\pandoc@box}%
  \ifdim\@tempb\p@<\@tempa\p@\let\@tempa\@tempb\fi% select the smaller of both
  \ifdim\@tempa\p@<\p@\scalebox{\@tempa}{\usebox\pandoc@box}%
  \else\usebox{\pandoc@box}%
  \fi%
}
% Set default figure placement to htbp
\def\fps@figure{htbp}
\makeatother
% definitions for citeproc citations
\NewDocumentCommand\citeproctext{}{}
\NewDocumentCommand\citeproc{mm}{%
  \begingroup\def\citeproctext{#2}\cite{#1}\endgroup}
\makeatletter
 % allow citations to break across lines
 \let\@cite@ofmt\@firstofone
 % avoid brackets around text for \cite:
 \def\@biblabel#1{}
 \def\@cite#1#2{{#1\if@tempswa , #2\fi}}
\makeatother
\newlength{\cslhangindent}
\setlength{\cslhangindent}{1.5em}
\newlength{\csllabelwidth}
\setlength{\csllabelwidth}{3em}
\newenvironment{CSLReferences}[2] % #1 hanging-indent, #2 entry-spacing
 {\begin{list}{}{%
  \setlength{\itemindent}{0pt}
  \setlength{\leftmargin}{0pt}
  \setlength{\parsep}{0pt}
  % turn on hanging indent if param 1 is 1
  \ifodd #1
   \setlength{\leftmargin}{\cslhangindent}
   \setlength{\itemindent}{-1\cslhangindent}
  \fi
  % set entry spacing
  \setlength{\itemsep}{#2\baselineskip}}}
 {\end{list}}
\usepackage{calc}
\newcommand{\CSLBlock}[1]{\hfill\break\parbox[t]{\linewidth}{\strut\ignorespaces#1\strut}}
\newcommand{\CSLLeftMargin}[1]{\parbox[t]{\csllabelwidth}{\strut#1\strut}}
\newcommand{\CSLRightInline}[1]{\parbox[t]{\linewidth - \csllabelwidth}{\strut#1\strut}}
\newcommand{\CSLIndent}[1]{\hspace{\cslhangindent}#1}

\makeatletter
\@ifpackageloaded{tcolorbox}{}{\usepackage[skins,breakable]{tcolorbox}}
\@ifpackageloaded{fontawesome5}{}{\usepackage{fontawesome5}}
\definecolor{quarto-callout-color}{HTML}{909090}
\definecolor{quarto-callout-note-color}{HTML}{0758E5}
\definecolor{quarto-callout-important-color}{HTML}{CC1914}
\definecolor{quarto-callout-warning-color}{HTML}{EB9113}
\definecolor{quarto-callout-tip-color}{HTML}{00A047}
\definecolor{quarto-callout-caution-color}{HTML}{FC5300}
\definecolor{quarto-callout-color-frame}{HTML}{acacac}
\definecolor{quarto-callout-note-color-frame}{HTML}{4582ec}
\definecolor{quarto-callout-important-color-frame}{HTML}{d9534f}
\definecolor{quarto-callout-warning-color-frame}{HTML}{f0ad4e}
\definecolor{quarto-callout-tip-color-frame}{HTML}{02b875}
\definecolor{quarto-callout-caution-color-frame}{HTML}{fd7e14}
\makeatother
\makeatletter
\@ifpackageloaded{caption}{}{\usepackage{caption}}
\AtBeginDocument{%
\ifdefined\contentsname
  \renewcommand*\contentsname{Table of contents}
\else
  \newcommand\contentsname{Table of contents}
\fi
\ifdefined\listfigurename
  \renewcommand*\listfigurename{List of Figures}
\else
  \newcommand\listfigurename{List of Figures}
\fi
\ifdefined\listtablename
  \renewcommand*\listtablename{List of Tables}
\else
  \newcommand\listtablename{List of Tables}
\fi
\ifdefined\figurename
  \renewcommand*\figurename{Figure}
\else
  \newcommand\figurename{Figure}
\fi
\ifdefined\tablename
  \renewcommand*\tablename{Table}
\else
  \newcommand\tablename{Table}
\fi
}
\@ifpackageloaded{float}{}{\usepackage{float}}
\floatstyle{ruled}
\@ifundefined{c@chapter}{\newfloat{codelisting}{h}{lop}}{\newfloat{codelisting}{h}{lop}[chapter]}
\floatname{codelisting}{Listing}
\newcommand*\listoflistings{\listof{codelisting}{List of Listings}}
\makeatother
\makeatletter
\makeatother
\makeatletter
\@ifpackageloaded{caption}{}{\usepackage{caption}}
\@ifpackageloaded{subcaption}{}{\usepackage{subcaption}}
\makeatother
\makeatletter
\@ifpackageloaded{tcolorbox}{}{\usepackage{tcolorbox}}
\makeatother
\makeatletter
\@ifpackageloaded{fontawesome5}{}{\usepackage{fontawesome5}}
\makeatother

\usepackage{bookmark}

\IfFileExists{xurl.sty}{\usepackage{xurl}}{} % add URL line breaks if available
\urlstyle{same} % disable monospaced font for URLs
\hypersetup{
  pdftitle={Protein modeling in the AlphaFold era},
  hidelinks,
  pdfcreator={LaTeX via pandoc}}


\title{Protein modeling in the AlphaFold era}
\author{}
\date{2023-08-21}

\begin{document}
\frontmatter
\maketitle

\renewcommand*\contentsname{Table of contents}
{
\setcounter{tocdepth}{1}
\tableofcontents
}

\mainmatter
\chapter{The history of protein structure modeling is told by a
competition
(CASP)}\label{the-history-of-protein-structure-modeling-is-told-by-a-competition-casp}

Every two years since 1994, groups in the field of structural
bioinformatics have conducted a worldwide experiment in which they
predict a set of unknown protein structures in a controlled, blind
test-like competition and compare their results with the experimentally
obtained structures. That is the \textbf{CASP} or
\href{https://predictioncenter.org/}{\emph{Critical assessment of
Protein Structure Prediction}}.

\begin{figure}

\centering{

\pandocbounded{\includegraphics[keepaspectratio]{pics/casp13.png}}

}

\caption{\label{fig-CASP13}Comparative Z-score of CASP13 participants.
The score is based in the GDT\_TS (Global distance test).}

\end{figure}%

The best research groups in the field test their new methods and
protocols in the CASP. However, at CASP13 (2018), an AI company called
\href{https://en.wikipedia.org/wiki/DeepMind}{Deepmind} (Google
Subsidiary) entered the scene. Their method, named Alphafold (Senior et
al. 2020) clearly won CASP13. Alphafold (v.1) implemented improvements
in some recently used approaches and created an entirely new pipeline.
Instead of creating contact maps from the alignment and then folding the
structure, they used an MRF unit (Markov Random Field) to extract the
main features of the sequence and MSA in advance and process all this
information into a multilayer NN (called ResNet) that also predicted
distance probabilities instead of contacts, resulting in high accuracy.
Then, Alphafold uses all the possibly obtained information to create the
structure and then improve it by energy minimization (steepest descent
method) and substitution of portions with a selected DB of protein
fragments.

\begin{figure}

\centering{

\pandocbounded{\includegraphics[keepaspectratio]{pics/alphafold1.png}}

}

\caption{\label{fig-alphafold1}Workflow of the first Alphafold method
presented in CASP13. MSA stands for multiple sequence alignment; PSSM
indicates Position-specific-scoring matrix and MRF stands for Markov
Random Field (or Potts model). From the Sergei Ovchinnikov and Martin
Steinegger presentation of Alphafold2 to the Boston Protein Design Group
(slides and video in the \hyperref[links]{link below}).}

\end{figure}%

After Alphafold, similar methods were also developed and made available
to the general public, like the \emph{trRosetta} (Yang et al. 2020),
from the \href{https://www.bakerlab.org/}{Baker lab}, available in
\href{https://www.rosettacommons.org/software}{Rosetta} open source
software and in the \href{https://robetta.bakerlab.org/}{Robetta}
server. This led to some controversy (mostly on Twitter) about the open
access to the CASP software and later on DeepMind publishes all the code
on GitHub.

\chapter{CASP14 or when protein structure prediction come to age for
(non structural)
biologists}\label{casp14-or-when-protein-structure-prediction-come-to-age-for-non-structural-biologists}

There was a lot of hype in CASP14 and the guys from DeepMind did not
disappoint anyone. Alphafold2 left all competitors far behind, both in
relative terms (score against the other groups) and in absolute terms
(lowest alpha-carbon RMSD). As highlighted earlier, the accuracy of many
of the predicted structures was \emph{within} the margin of error of the
experimental determination methods (see for instance Mirdita et al.
2022).

\begin{figure}

\centering{

\href{https://predictioncenter.org/casp14/zscores_final.cgi?formula=gdt_ts}{\pandocbounded{\includegraphics[keepaspectratio]{pics/casp14.png}}}

}

\caption{\label{fig-casp14}Comparative CASP14 scores}

\end{figure}%

\begin{figure}

\centering{

\href{https://www.nature.com/articles/s41586-021-03819-2/figures/1}{\pandocbounded{\includegraphics[keepaspectratio]{pics/jumper2021.png}}}

}

\caption{\label{fig-alphafold2bench}Performance of Alphafold2 the CASP14
dataset relative to the top-15 entries. Data are median and the 95\%
confidence interval of the median, for alpha-carbom RMSD. Panels b-c-d
show example comparison between model and experimental structures.}

\end{figure}%

Deepming took some time (eight months, which is an eternity nowadays) to
publish the method (Jumper et al. (2021)) and to release the code on
\href{https://github.com/deepmind/alphafold}{Github}, but other new
methods, such as RoseTTAfold (Baek et al. (2021)) and C-I-Tasser (Zheng
et al. (2021)) were able to obtain similar results and were available on
public servers, which may have pushed Deepmind to make it all available
to the scientific community\footnote{See
  https://www.wired.com/story/without-code-for-deepminds-protein-ai-this-lab-wrote-its-own/}.
Not surprisingly, a group of independent scientists
(\href{https://twitter.com/sokrypton}{Sergey Ovchinnikov},
\href{https://twitter.com/milot_mirdita}{Milot Mirdita}, and
\href{https://twitter.com/thesteinegger}{Martin Steinegger}), decided to
implement Alphafold2 in a
\href{https://github.com/sokrypton/ColabFold}{Colab notebook}, called
\textbf{ColabFold} Mirdita et al. (2022), which freely available online
using \emph{Google Colab} notebooks platform. Other free implementations
of Alphafold have been and are available, but ColabFold has been the
most widely discussed and well known. They implemented some tricks to
speed up the modeling, most notably the use of
\href{https://docs.google.com/presentation/d/1mnffk23ev2QMDzGZ5w1skXEadTe54l8-Uei6ACce8eI/present\#slide=id.ge58c80b745_0_15}{MMSeqs2}
(developed by Martin Steinegger's group) to search for homologous
structures on Uniref30, which made Colabfold a fast method that made all
the previous advanced methods almost useless. This was the real
breakthrough in the field of protein structure prediction, making
Alphafold accessible and, also very importantly, facilitated the further
development of the method, implementing very quickly new features, like
the prediction of protein complexes, which was actually first mentioned
on
\href{https://docs.google.com/presentation/d/1mnffk23ev2QMDzGZ5w1skXEadTe54l8-Uei6ACce8eI/present\#slide=id.ge58c80b745_0_15}{Twitter}
and then led to several new dedicated methods, including
AlphaFold-multimer (Evans et al. 2022) from Deepmind or virtual PullDown
approaches (Yu et al. 2023) (Shin et al. 2025).

\chapter{Alphafold as the paradigm of a New Era}\label{sec-AF2}

\section{Why is Alphafold so freaking
accurate?}\label{why-is-alphafold-so-freaking-accurate}

The emergence of the AlphaFold era is the result of decades of previous
work that progressively incorporated neural networks (NNs) and deep
learning into protein structure prediction. Because NNs can learn subtle
patterns and correlations from very large datasets, they are
exceptionally well suited to the challenge of inferring 3D protein
structures from sequence information. Over the last 20--30 years,
advances in protein structure prediction have closely followed
developments in deep learning, gradually integrating convolutional
neural networks (CNNs), recurrent neural networks (RNNs), and eventually
Transformer models (Kumar et al. 2025).

The Transformer architecture, originally designed for machine
translation, achieved outstanding performance thanks to its
encoder--decoder structure. The encoder converts the input sequence into
internal vector representations, while the decoder generates the output.
Its key innovation is the multi-head attention mechanism, which
identifies complex relationships between tokens by assigning attention
weights according to their relevance. This enables the model to capture
long-range dependencies in sequences far more effectively than previous
architectures. Transformers typically consist of six encoder layers and
six decoder layers, each including multi-head attention and feed‑forward
networks. These components are highly parallelizable, which makes
Transformer-based models both efficient and powerful (Wang, Peng, and
Yang 2022).

The core idea behind AlphaFold v2 (referred to here simply as AlphaFold)
and related methods is analogous to strategies used in image processing.
In computer vision, a 2D image is represented as a 3D tensor because of
the RGB or HSV color channels. Protein structure prediction uses a
similar approach: 1D and 2D sequence-derived features are transformed
into a 3D volume with multiple channels encoding inter-residue
relationships. This rich, multi-channel tensor serves as the input to
deep learning models that predict distances and structural constraints
between residues (Pakhrin et al. 2021).

\begin{figure}

\centering{

\pandocbounded{\includegraphics[keepaspectratio]{pics/machine_fold.png}}

}

\caption{\label{fig-deep}From the perspective of Deep Learning method
development, the problem of protein distogram or real-valued distance
prediction (bottom row) is similar to the `depth prediction problem' in
computer vision. From Pakhrin et al. (2021).}

\end{figure}%

Alphafold can be explained as a pipeline with three interconected tasks
(see picture below). First, in contrast to Alphafold v.1, the input to
Alphafold is a ``raw'' MSA, i.e., the deep learning network extracts the
co-evolutionary information directly from the MSA. It queries several
databases of protein sequences and constructs an MSA that is used to
select templates. This can be a limiting step, affecting the speed of
modeling (see below), but it can be also related to model accuracy, as
has been recently shown in CASP15 (Lee et al., n.d.; Peng et al. 2023).

In the second part of the diagram, AlphaFold takes the multiple sequence
alignment and the templates, and processes them in a Transformer that
they named \emph{Evoformer}. This process has been referred by some
authors as inter-residue interaction map-threading (Bhattacharya et al.
2021). The objective of this part is to extract layers of information to
generate residue interaction maps. The Evoformer refines both its
understanding of the MSA and the predicted 3D structure in an iterative
process to refine the model. Improved predictions of amino acid
interactions help refine the MSA representation, and vice-versa.

The AF2 transformer architecture allows the model to process all amino
acids in a protein sequence simultaneously, rather than one at a time,
through a technique known as the attention mechanism. This is a
significant advantage because it enables the model to capture long-range
relationships between amino acids. Essentially, the attention mechanism
helps a neural network to concentrate on the most relevant parts of its
input when making predictions. For proteins, this means the network can
learn which amino acids are likely to interact with each other, even if
they are far apart in the sequence. It's akin to reading a lengthy
document; you don't focus equally on every word but rather on the key
phrases that are crucial for understanding the overall meaning.

Importantly, in the AF2 Evoformer, this process is iterative and the
information goes back and forth throughout the network. At each
recycling step, the model refines its predictions of amino acid
relationships, leading to a more accurate structure (the original model
uses 3 cycles). As explained in the great
\href{https://www.blopig.com/blog/2021/07/alphafold-2-is-here-whats-behind-the-structure-prediction-miracle/}{post}
from Carlos Outerial at the OPIG site:

\begin{quote}
\emph{This is easier to understand as an example. Suppose that you look
at the multiple sequence alignment and notice a correlation between a
pair of amino acids. Let's call them A and B. You hypothesize that A and
B are close, and translate this assumption into your model of the
structure. Subsequently, you examine said model and observe that, since
A and B are close, there is a good chance that C and D should be close.
This leads to another hypothesis, based on the structure, which can be
confirmed by searching for correlations between C and D in the MSA. By
repeating this several times, you can build a pretty good understanding
of the structure.}
\end{quote}

The third part of the pipeline is the structure building module, which
uses the information from the previous steps to construct a 3D model
structure protein of the query sequence. This network will give you a
single model, without heavy energy optimization steps (only a relax
protocol). Model building is based in a new concept of 3D structures
generation, named \textbf{IPA} (Invariant Point Attention), a core
neural network module that uses of a curated list of parametrized list
of torsion angles to generate the side chains. IPA allows the network to
maintain rotational and translational invariance, which is crucial for
accurate 3D structure prediction. Earlier attempts to develop and
end-to-end method were unsuccessful because the structure representation
was not optimal. Even methods implemented after AlphaFold, like
RoseTTAFold, use less efficient methods and often predict very quickly
and accurately the backbone coordinates but require external programs to
generate an all-atoms model.

\begin{figure}

\centering{

\pandocbounded{\includegraphics[keepaspectratio]{pics/alphafols2.png}}

}

\caption{\label{fig-af2}Oxford Proteins Informatics Group Blog, modified
from Jumper et al. (2021)}

\end{figure}%

Like for most of the previous methods Alphafold would give your better
results with proteins with related structures known and with a lot of
homologs in Uniref databases. However, comparing to nothing, it will
likely give you (limited) useful results for the so-called ``dark
genome''. I work with phages and bacterial mobile elements, and
sequencing that is often frustrating as more than 50\% of the proteins
have no homologous in the database. So you have a bunch of proteins of
unknown function\ldots{} However, as we do know that structure is more
conserved than sequence, we may use the structure to find out the
function of our dark proteins. There are a few resources for this, I'd
suggest you to try \href{https://search.foldseek.com/search}{FoldSeek}
(M. van Kempen et al., n.d.) and
\href{http://ekhidna2.biocenter.helsinki.fi/dali/}{Dali} (Holm 2022)
servers. You can upload the PDB file of your model and search for
related structures in RCSB PDB database and also in Alphafold database.

\begin{tcolorbox}[enhanced jigsaw, toptitle=1mm, colbacktitle=quarto-callout-tip-color!10!white, breakable, opacityback=0, coltitle=black, arc=.35mm, colback=white, leftrule=.75mm, titlerule=0mm, colframe=quarto-callout-tip-color-frame, opacitybacktitle=0.6, bottomtitle=1mm, title=\textcolor{quarto-callout-tip-color}{\faLightbulb}\hspace{0.5em}{Tip}, bottomrule=.15mm, toprule=.15mm, left=2mm, rightrule=.15mm]

\emph{FoldSeek} needs only a few seconds/minutes and is therefore faster
than \emph{Dali}. Therefore, it is better to use Dali for some selected
searches that require a double check or a more reliable result, even if
it may take a few days.

\end{tcolorbox}

As mentioned above, Colabfold aims to make the process faster by using
MMSeqs in the first block. Additionally, the number of recycling steps
can also be adapted. Moreover, different Colabfold notebooks have been
developed (and evolved) to allow some customization and other feature,
like batch processing of multiple proteins avoiding recompilation and
identification of protein-protein interactions (Mirdita et al. 2022).

Alphafold models can be evaluated by the mean \textbf{pLDDT}, a
per-residue confidence metric. It is stored in the B-factor fields of
the mmCIF and PDB files available for download (although unlike a
B-factor, higher pLDDT is better). The model confidence can vary greatly
along a chain so it is important to consult the confidence when
interpreting structural features. Very often, the lower confidence
fragments are not product of a poor prediction but an indicator of
protein disorder (Wilson, Choy, and Karttunen 2022). Additionally, AF2
calculates the predicted template modeling, or TM-scores (\textbf{pTM})
based on a pairwise error prediction, and is used to calculate the
predicted aligned error (\textbf{PAE}) that estimates the error of the
position of each amino acid. Interestingly, the PAE can be used as a
predictor for protein dynamics, as PAE correlates with distance
variations from the MD simulations (Guo et al. 2022).

Alphafold also partnered with EMBL-EBI and Uniprot and generated a huge
curated database of proteins from model organisms (Varadi et al. 2022),
the \href{https://alphafold.ebi.ac.uk/}{Alphafold database}. This is an
amazing resource that may be also very helpful for you. Just consider
that this database increased from 48\% to 76\% the fraction of human
proteome with structural data, and also it also means great increases in
the case of other model organisms, like, including microorganisms and
plants (Porta-Pardo et al. 2022).

\begin{figure}

\centering{

\includegraphics[width=4.51042in,height=\textheight,keepaspectratio]{pics/journal.png}

}

\caption{\label{fig-afDDBB}Changes in protein structural coverage in
model organisms (From Porta-Pardo et al. (2022)).}

\end{figure}%

\section{\texorpdfstring{{Wanna try
Alphafold?}}{Wanna try Alphafold?}}\label{wanna-try-alphafold}

\emph{Section under construction!}

As mentioned above, the grand breakthrough of Alphafold would not have
been the same without the Colabfold, a free open tool that made the
state-of-the-art of AI-fueled protein prediction available to everyone.

\begin{figure}

\href{https://github.com/sokrypton/ColabFold}{\includegraphics[width=4.59375in,height=\textheight,keepaspectratio]{pics/qrcode.png}}

\caption{\label{fig-colabfold}ColabFold GitHub repository}

\end{figure}%

The \href{https://github.com/sokrypton/ColabFold}{Colabfold} repository
on GitHub contains links to several Python notebooks developed on
\href{https://colab.research.google.com/}{Google Colab}, a platform to
develop and share Python scripts on a Jupyter Notebook format. Notebooks
are very important also for reproducibility in computer sciences, as
they allow you to have the background and details and the actual code in
a single document and also execute it. You can share those notebooks
very easily and also update quickly as they are stored in your Google
Drive.

Colabfold allow you to run notebooks of Alphafold, RoseTTAfold, and
ESMFold for specific applications, allowing even to run a bunch of
proteins in batch. You can see a more detailed description in Mirdita et
al. (2022). We are using the
\href{https://colab.research.google.com/github/sokrypton/ColabFold/blob/main/AlphaFold2.ipynb}{Alphafold2\_mmseqs2}
notebook, that allow you most of the common features. You need to allow
Colab to use your Google account.

\begin{figure}

\centering{

\pandocbounded{\includegraphics[keepaspectratio]{pics/colab1.jpg}}

}

\caption{\label{fig-colabfold2}Introducing your sequence in Colabfold}

\end{figure}%

Then paste your sequence and chose a name. For more accurate models you
can click ``use\_amber'' option. It will run a short
\href{https://en.wikipedia.org/wiki/Molecular_dynamics}{\emph{Molecular
Dynamics}} protocol that ultimately optimize the modeling, but it will
also take some more time, so better try at home.

As you can see, an this is a recent feature, you can also add your own
template, which combines the classical ``homology modeling'' based upon
a single template with the advanced AlphaFold model construction method.
That will safe time, but of course without any guarantee. If you have a
template of a related protein, like an alternative splicing or a disease
mutant, I'd advise you to try with and without the template. You may
surprise.

\begin{figure}

\centering{

\pandocbounded{\includegraphics[keepaspectratio]{pics/colab_execute.jpg}}

}

\caption{\label{fig-colabfold3}Executing Colabfold}

\end{figure}%

At this point, you may execute the whole pipeline or may some more
customization. MSA stage can be also optimized to reduce execution time,
by reducing database or even by providing your own MSA. Very often you
may want to fold a protein with different parameters, which may very
convenient to reuse an MSA from a previous run (although they recently
updated servers for MMSeqs and made it really faster). If your proteins
are in the same operon or by any other reason you think that they should
have co-evolved, you prefer a ``paired'' alignment. But you can always
do both.

Advanced settings are specially needed for protein-protein complexes.
Also the number of recycling steps will improve your model, particularly
for targets with no MSA info from used databases. Then you can just get
your model (and companion info and plots) in your GDrive or download it.

{\textbf{What do you think is the ideal protein for AlphaFold? Do you
think homology modeling is dead?}}

\section{Now what? The Post-AlphaFold era
(2022-2023)}\label{now-what-the-post-alphafold-era-2022-2023}

As mentioned earlier, RoseTTAFold was released at the same time as
AlphaFold's paper and code, although it is clearly inspired by
AlphaFold's capabilities in CASP14. It is based on a three-track
network, and recent implementations have allowed prediction of protein
interactions with nucleic acids. Other methods such as AlphaFold and
RoseTTAFold were released later, as were
\href{https://openfold.io/}{OpenFold} and
\href{https://colab.research.google.com/github/dptech-corp/Uni-Fold/blob/main/notebooks/unifold.ipynb}{UniFold},
which are based on the PyTorch Transformers AI framework. More recently,
RoseTTAFold2, which extends the three-track architecture over the whole
network and incorporate other new advances and AlphaFold tricks, like
the FAPE (frame aligned point error) or recycling steps during the
training, giving rise to an end-to-end model equivalent in accuracy to
AF2 for monomers and AF2-multimer for complexes, with better
computational scaling on proteins and complexes larger than 1000
residues (Baek et al., n.d.).

The use of MSA has been cited as a limitation for AlphaFold and related
methods. However, the predictions are significantly worse without MSA or
with MSAs without depth. One way to improve predictions is to use a
protein natural language model (\textbf{PLMs}), i.e., a model trained to
predict sequence from sequence and not dependent on good MSAs.
\href{https://github.com/HeliXonProtein/OmegaFold}{Omegafold} and
\href{https://esmatlas.com/resources?action=fold}{ESMfold} (also in
Colabfold,
\href{https://colab.research.google.com/github/sokrypton/ColabFold/blob/main/beta/omegafold.ipynb}{here}
and
\href{https://colab.research.google.com/github/sokrypton/ColabFold/blob/main/ESMFold.ipynb}{here},
respectively) are two new implementations that require only a single
sequence. While AlphaFold2 uses deep learning with multiple sequence
alignments (MSAs) to infer protein structures, PLMs learn the
``language'' of protein sequences directly, representing amino acid
sequences. Each amino acid (e.g., alanine, glycine) is represented as a
``token,'' a numerical representation that the model can process. By
training on extensive datasets of protein sequences, PLMs learn the
patterns and relationships between amino acids, enabling them to predict
3D structures from single sequences. Essentially, PLMs treat amino acid
sequences as ``sentences,'' where each amino acid is a ``word''. This
method offers advantages in speed and applicability, particularly for
proteins without sufficient sequence homology. However, they may have
lower overall accuracy compared to AlphaFold2 when MSAs are available.
PLMs perform efficiently when using a single sequence, often surpassing
AlphaFold in this scenario. Nevertheless, although they can outperform
AF2, since PLMs are trained on existing sequences, these methods still
show reduced performance on orphan sequences. That is, the language
models seem to replicate the training MSAs (Elofsson 2023).

The appearance of ESMfold and the companion
\href{https://esmatlas.com/}{Metagenomic Atlas} (Lin et al. 2023) was
seen as a new step that could trigger a new revolution in the field,
since it was developed by scientists at META (Callaway 2022). So it was
now a sort of Google versus Facebook battle for the most powerful AI
methods for protein modeling. Last summer, however, we learned that META
had decided to discontinue the ESM project.

\begin{centering}\begin{tcolorbox}[hbox,
colframe=lightgray,
colback=white]

\textcolor{ProcessBlue}{{\Large {\faTwitterSquare}}} %

\href{https://twitter.com/KevinKaichuang/status/1689071501998260224}{\raisebox{0.1 em}{Click to view embedded Twitter post.}}

\end{tcolorbox}\end{centering}
\vspace{1em}

AlphaFold was again the protagonist of CASP15, which took place in 2020,
but it did not overwhelm the rest like in 2018.

\begin{centering}\begin{tcolorbox}[hbox,
colframe=lightgray,
colback=white]

\textcolor{ProcessBlue}{{\Large {\faTwitterSquare}}} %

\href{https://twitter.com/sokrypton/status/1601596893355642882}{\raisebox{0.1 em}{Click to view embedded Twitter post.}}

\end{tcolorbox}\end{centering}
\vspace{1em}

AlphaFold was used in some form in more than half of the protocols, and
while the standard AlphaFold method performed better than many other
methods, several groups achieved significant improvements for monomeric
proteins and protein assemblies. In short (see Elofsson (2023)), we
learned at CASP15 that there are two main ways to improve AlphaFold: (1)
more efficient use of templates, increasing the size of the database or
sampling through more efficient MSAs, or (2) hacking AlphaFold to use
dropouts to generate thousands of models for each target, which
increases computational time but also increases the chances of better
models. Other small improvements have been proposed, such as refinement
steps (Adiyaman et al. 2023) or using improved template search or new
scoring capabilities based on Voronoi surfaces or Deep Learning.

\subsection{\texorpdfstring{2021-2023 - Short summary in one
wire}{2021-2023 - Short summary in one  wire}}\label{short-summary-in-one-wire}

\begin{centering}\begin{tcolorbox}[hbox,
colframe=lightgray,
colback=white]

\textcolor{ProcessBlue}{{\Large {\faTwitterSquare}}} %

\href{https://twitter.com/simocristea/status/1626304239931912192}{\raisebox{0.1 em}{Click to view embedded Twitter post.}}

\end{tcolorbox}\end{centering}
\vspace{1em}

\section{Corollary 1: Did Deepmind ``fold'' Levinthal's
paradox?}\label{corollary-1-did-deepmind-fold-levinthals-paradox}

The development of Alphafold and the
\href{https://alphafold.ebi.ac.uk/}{Alphafold structures Database} in
collaboration with \href{https://www.ebi.ac.uk/about}{EMBL-EBI} was been
the origin of a New Era. Moreover, in a further turn, the
\href{https://esmatlas.com/}{Metagenomic Atlas} by Meta AI uncovers
virtually the whole protein space. Thanks to these milestones the
coverage of the protein structure space has been greatly increased
(Porta-Pardo et al. 2022), which virtually close the sequence-structure
gap. Since 2020, many scientific publications and journals worldwide
published long articles about the meaning of this breakthrough in
science and its applications in biotechnology and biomedicine\footnote{https://www.bbc.com/news/science-environment-57929095

  https://www.forbes.com/sites/robtoews/2021/10/03/alphafold-is-the-most-important-achievement-in-ai-ever/

  https://elpais.com/ciencia/2021-07-22/la-forma-de-los-ladrillos-basicos-de-la-vida-abre-una-nueva-era-en-la-ciencia.html}
and DeepMind claimed to have
\href{https://www.deepmind.com/blog/alphafold-a-solution-to-a-50-year-old-grand-challenge-in-biology}{Solved
a 50-years Grand Challenge in biochemistry}.

In other words, after AlphaFold, would it no longer be necessary to
perform X-ray crystallography or nuclear magnetic resonance? First,
AlphaFold models can be used in electron density maps and help solve
complex cases. Thus, the new framework helps crystallographers focus
their work on the most unresolved and difficult structures, such as
coiled-coils or holoproteins, which cause modulations and challenges in
the development of crystallographic methods.

However, some scientists argued that Alphafold2 and RoseTTAfold actually
\emph{cheat} as they do not really solve the problem but generate a deep
learning pipeline that is able to bypass the problem (Pederson 2021). In
agreement with that, it has been shown that machine learning methods
actually do not reproduce the expected folding pathways while improving
the structures during the recycling steps (Outeiral, Nissley, and Deane,
n.d.).

In conclusion, I believe that the Levinthal paradox is not (yet) fully
solved, although it seems to be close (Al-Janabi 2022). Practically, it
is solved for most of the protein space, but if your protein does not
have a homolog in the databases, you will still have some open
questions.

\chapter{Protein hallucination and diffusion
models}\label{protein-hallucination-and-diffusion-models}

Biochemists and Chemists have dreamed of designing proteins with
customizable properties for decades. First successful protein designs
were published in the late 1990s and several breakthrough by Baker's lab
pushed forward the idea. More recently, artificial intelligence is
starting a revolution in this field, and it will radically change it
forever (reviewed in Notin et al. (2024)).

Soon after AF2 and RoseTTAFold publication, they were used to
``hallucinate'' new proteins. In December 2021, Baker and his colleagues
reported expressing 129 of these hallucinated proteins in bacteria, and
found that about one-fifth of them folded into something resembling
their predicted shape. Then, in 2023, they introduced
\href{https://www.bakerlab.org/2023/03/30/rf-diffusion-now-free-and-open-source/}{RFdiffussion},
a new paradigm in protein design, based in AI diffusion models. Imagine
starting with a cloud of randomly positioned atoms (pure noise). During
training, diffusion models apply usually-random noise to a distribution
of data until it resembles pure static, for example. The model then
learns to reverse this process, generating final data points
representative of the original data distribution. This process is guided
by learned patterns of how atoms and molecules interact, allowing the
model to ``denoise'' the initial random arrangement into a realistic
biomolecular complex. The model learns to reverse a ``diffusion
process,'' where noise is progressively added to known biomolecular
structures. By learning to undo this process, the model can generate new
structures from scratch.

More recently several groups have combined diffusion models with NLPs to
develop improved methods to design new proteins. Protein sequences can
be described as a concatenation of letters from a chemically defined
alphabet, the natural amino acids, and like human languages, these
letters arrange to form secondary structural elements (``words''), which
assemble to form domains (``sentences'') that undertake a function
(``meaning''). One of the most attractive similarities is that protein
sequences, like natural languages, are information-complete: they store
structure and function entirely in their amino acid order with extreme
efficiency. With the extraordinary advances in the NLP field in
understanding and generating language with near-human capabilities, we
hypothesized that these methods open a new door to approach
protein-related problems from sequence alone, such as protein design.

You can try some of these state-of-the-art method using Colab notebooks
from \href{https://github.com/sokrypton/ColabDesign}{ColabDesing}
project, again by Sergey Ovchinnikov and colleages (check out his talks
\href{https://www.youtube.com/watch?v=2HmXwlKWMVs}{here} and
\href{https://www.youtube.com/watch?v=smjYJSudwr0}{here}).

\begin{centering}\begin{tcolorbox}[hbox,
colframe=lightgray,
colback=white]

\textcolor{ProcessBlue}{{\Large {\faTwitterSquare}}} %

\href{https://twitter.com/sokrypton/status/1514227001183416320}{\raisebox{0.1 em}{Click to view embedded Twitter post.}}

\end{tcolorbox}\end{centering}
\vspace{1em}

\chapter{Protein structural
phylogenetics}\label{protein-structural-phylogenetics}

Phylogenetic inference can be performed using either DNA or protein
sequences. DNA‑based phylogenetics is often preferable when comparing
closely related genes or species. However, as evolutionary distances
increase, nucleotide substitution matrices tend to saturate, making
protein‑based analyses more informative. Even so, deep sequence‑based
phylogenetics remains challenging. A major obstacle is substitution
saturation, where individual sites in an alignment accumulate multiple
substitutions along long branches. When this occurs, and especially when
substitution models fail to fully capture the true evolutionary process,
saturation can generate misleading signals and produce artifacts in
phylogenetic trees. Structural phylogenetics offers a promising
alternative to overcome these limitations. Protein structures evolve
much more slowly than their underlying sequences, and if models of
structural evolution can be inferred, they may help resolve phylogenetic
relationships that lie beyond the reach of classical sequence‑based
methods. The development of protein language alphabets that encode
structural information now enables the integration of these ``structural
characters'' within standard amino‑acid likelihood frameworks. Examples
of such approaches include FoldTree (Moi et al. 2025) and 3DiPhy
(Puente-Lelievre et al., n.d.), both of which use the set of 20 3Di
structural characters originally introduced in FoldSeek (M. van Kempen
et al. 2023). More recently, new 3Di‑based substitution matrices have
further improved the performance of these structural phylogenetic
methods (Garg and Hochberg 2025).

\chapter{Alphafold3 expanded the
scope}\label{alphafold3-expanded-the-scope}

Proteins are composed of amino acids, but they also include other
components.
\href{https://en.wikipedia.org/wiki/Holoprotein}{Holoproteins} often
contain multiple protein chains, ions, and cofactors. Additionally,
proteins' functions typically involve interactions with various ligands,
including DNA or RNA molecules. One limitation of AlphaFold 2 is its
ability to predict \emph{only} the structure of amino acid chains. It
generally predicts the orientation of amino acids to some binding
pockets accurately, particularly for ions or other common small ligands.
Furthermore, it has been demonstrated from the beginning that AF2 can
predict protein-protein interactions effectively, which was enhanced
with the Alphafold-Multimer version (Evans et al. 2022) and further
optimized in methods allowing full interactome predictions (Yu et al.
2023; Molodenskiy et al., n.d.; Rouger et al., n.d.). However, many
applications of protein structures in biology, biotechnology, and
biomedicine require the interaction of protein chains with other
molecules, such as organic ligands, drugs, or DNA/RNA molecules,
presenting new challenges in the field of protein prediction.

AlphaFold3 can now predict the structures and interactions of a much
wider range of biomolecules, including protein-protein complexes,
protein-DNA interactions, protein-RNA interactions, and small
molecule-protein interactions. This leap in capability is powered by two
key innovations: diffusion models and tokenization. DeepMind proposed
that AlphaFold3 has the potential to revolutionize drug discovery,
materials science, and our understanding of fundamental biological
processes. By accurately modeling the intricate interactions of
biomolecules, it can accelerate the development of new therapies,
materials, and technologies. For instance, AlphaFold3 could help predict
how a person's specific protein structures and mutations will respond to
different drugs, leading to highly effective, personalized treatments.
This could revolutionize healthcare, making treatments more precise and
reducing the trial-and-error approach currently used. On a global scale,
AF3's predictions could aid in rapidly developing vaccines and
treatments for emerging diseases, helping to prevent pandemics and
improve global health security. Its impact on neglected diseases could
be significant, bringing new hope to millions in underdeveloped regions.
This could lead to a more equitable distribution of healthcare
advancements, improving quality of life worldwide.

Imagine AlphaFold as the stage of learning to draw a single shape, and
AlphaFold2 as mastering the perfection of that shape. AlphaFold3, on the
other hand, is like learning to draw that shape and understanding how it
fits together with various other shapes to create an intricate and
complex picture.

\section{Is AF3 that different?}\label{is-af3-that-different}

Unlike AlphaFold2's direct prediction approach, AlphaFold3 uses a
diffusion model, a technique borrowed from image generation. The
diffusion process allows the model to consider the overall context of
the system, leading to more accurate predictions of complex
interactions.

AlphaFold3 (Abramson et al. 2024a) treats biomolecules as ``tokens,''
similar to how words are represented in natural language processing.
This allows the model to apply the same powerful language modeling
techniques used in large language models (LLMs). Instead of amino acid
sequences, the model works with a broader vocabulary of tokens
representing atoms, residues, nucleotides, and small molecules. Each
type of atom, residue, or molecule is assigned a unique token. The model
learns the statistical relationships between these tokens, capturing the
rules of biomolecular interactions. This enables AlphaFold3 to
``understand'' the language of biomolecules, predicting how they will
interact and assemble. Tokenization provides a unified representation
for diverse biomolecules, enabling the model to handle complex systems.
This, in combination with transformers, allows the model to understand
the context of biomolecular interactions, capturing long-range
dependencies and intricate relationships.

\begin{tcolorbox}[enhanced jigsaw, toptitle=1mm, colbacktitle=quarto-callout-warning-color!10!white, breakable, opacityback=0, coltitle=black, arc=.35mm, colback=white, leftrule=.75mm, titlerule=0mm, colframe=quarto-callout-warning-color-frame, opacitybacktitle=0.6, bottomtitle=1mm, title=\textcolor{quarto-callout-warning-color}{\faExclamationTriangle}\hspace{0.5em}{Warning}, bottomrule=.15mm, toprule=.15mm, left=2mm, rightrule=.15mm]

It is important to note that while AlphaFold3 incorporates the concept
of tokenization from language models, it utilizes it within a unique
framework for a broader purpose. Tokenization in AlphaFold3 is used to
represent all components of a complex, enabling the diffusion model to
process these components simultaneously. Unlike protein language models,
which primarily focus on learning statistical patterns from large
training sets of protein sequences, AlphaFold3 relies on the diffusion
model facilitated by tokenization.

\textbf{Therefore, it is more accurate to describe AlphaFold3 as a
\emph{hybrid} or \emph{evolved} model that integrates core principles of
language modeling rather than integrating (at any step) protein language
models.}

\end{tcolorbox}

Unlike its predecessor, AF3 is a generative model. This means it can
give different outputs for the same input and that it can hallucinate, a
phenomenon evident with intrinsically disordered regions (IDRs) of
proteins.

AF3 can be tested in a Deepmind's
\href{https://alphafoldserver.com/}{server} that allows protein:protein,
protein:DNA/RNA interactions, as well as interactions with a small
subset of molecules.

\section{Did AF3 move forward the field like
AF2?}\label{did-af3-move-forward-the-field-like-af2}

Short answer: Yes.

There are two main reasons. First, once they published the paper, the
community moved to analyze the AF3 workflow and innovations and many
labs attempted to reproduce or even improve it. Additionally, like in
the case of AF2, Deepmind did not initially release the full code. In
this case, they even published the Nature paper without the code. That
generated a huge controversy in the social media (again):

\begin{centering}\begin{tcolorbox}[hbox,
colframe=lightgray,
colback=white]

\textcolor{ProcessBlue}{{\Large {\faTwitterSquare}}} %

\href{https://twitter.com/simocristea/status/1788515833884389806}{\raisebox{0.1 em}{Click to view embedded Twitter post.}}

\end{tcolorbox}\end{centering}
\vspace{1em}

\begin{centering}\begin{tcolorbox}[hbox,
colframe=lightgray,
colback=white]

\textcolor{ProcessBlue}{{\Large {\faTwitterSquare}}} %

\href{https://twitter.com/simocristea/status/1788515833884389806}{\raisebox{0.1 em}{Click to view embedded Twitter post.}}

\end{tcolorbox}\end{centering}
\vspace{1em}

There was also some discussion with comments and editorials in Nature
and other journals ({``AlphaFold3 {\textemdash} Why Did Nature Publish
It Without Its Code?''} 2024; Callaway 2024) and a public letter with
more than 100 signatures (Wankowicz et al. 2024). Ultimately, after the
Nobel prize concession, they decided to publish an addendum to the paper
and release the code on github (Abramson et al. 2024b).

Currently, besides the AF3 released code, there are several AF3-based or
AF3 alternatives with similar performance that can be tested by any
user. Some examples:

\begin{enumerate}
\def\labelenumi{\arabic{enumi}.}
\item
  \href{https://github.com/bytedance/Protenix}{Protenix} is a fully
  opensource AF3 implementation based on PyTorch that allow handling any
  molecule (Team et al., n.d.)
\item
  Chai-1 (Discovery et al., n.d.) is similar to AF3 and it can be used
  as a Python package for non-commercial use and via a
  \href{https://lab.chaidiscovery.com/}{web interface} where it can be
  used for free including for commercial drug discovery purposes. It
  uses
  \href{https://en.wikipedia.org/wiki/Simplified_Molecular_Input_Line_Entry_System}{smiles}
  encoding for representation of any ligand molecule
\item
  \href{https://github.com/jwohlwend/boltz}{Boltz-1} (Wohlwend et al.,
  n.d.) is also a very recent method that achieves AF3-level accuracy in
  predicting the 3D structures of biomolecular complexes and it can be
  tested in Colabfold.
\end{enumerate}

Finally, I recently came into
\href{https://github.com/colbyford/PyMOLfold}{PyMOLFold} (Ford et al.,
n.d.), a novel Pymol plugin that facilitate the use of these methods. It
provides instructions for installation (on different Conda environments)
and they work on your computer. I tested it using only CPUs and it
worked, though quite slowly as compared with the use of online servers
or Colabs.

\begin{centering}\begin{tcolorbox}[hbox,
colframe=lightgray,
colback=white]

\textcolor{ProcessBlue}{{\Large {\faTwitterSquare}}} %

\href{https://twitter.com/c01_by/status/1867227053197574556}{\raisebox{0.1 em}{Click to view embedded Twitter post.}}

\end{tcolorbox}\end{centering}
\vspace{1em}

\section{Corollary 2: A Nobel prize to the
field}\label{corollary-2-a-nobel-prize-to-the-field}

Between 2021 and 2024, I used to conclude my lectures on protein
modeling by proposing a wager to my students, predicting that the
creators of AF2 and RoseTTAFold would win a Nobel Prize. Indeed, the
\href{https://www.nobelprize.org/prizes/chemistry/2024/press-release/}{Nobel
Prize in Chemistry for 2024} was awarded to \textbf{David Baker} (1/2)
for ``computational protein design'' and to \textbf{Demis Hassabis}
(1/4) along with \textbf{John Jumper} (1/4) for ``protein structure
prediction''.

As stated by David Baker (Baker and Lin 2025), this Nobel Prize serves
as recognition for the field and an encouragement for future
advancements, highlighting the significant impact that protein design
can have on the world.

\chapter{\texorpdfstring{{Review Quizz}}{Review Quizz}}\label{quizz}

\textbf{AlphaFold and RoseTTAFold are two programs that have been
designed to?}

\begin{itemize}
\tightlist
\item
  \begin{enumerate}
  \def\labelenumi{(\Alph{enumi})}
  \tightlist
  \item
    Predict how proteins fold (The protein folding problem)\\
  \end{enumerate}
\item
  \begin{enumerate}
  \def\labelenumi{(\Alph{enumi})}
  \setcounter{enumi}{1}
  \tightlist
  \item
    Predict putative binding sites on proteins (for drug design)\\
  \end{enumerate}
\item
  \begin{enumerate}
  \def\labelenumi{(\Alph{enumi})}
  \setcounter{enumi}{2}
  \tightlist
  \item
    Predict the native structure of a protein (The protein structure
    prediction problem)\\
  \end{enumerate}
\item
  \begin{enumerate}
  \def\labelenumi{(\Alph{enumi})}
  \setcounter{enumi}{3}
  \tightlist
  \item
    Predict if two proteins interact, and, if they do, the structure of
    the complex
  \end{enumerate}
\end{itemize}

\textbf{Which of these statements related to AlphaFold is most likely to
be incorrect?}

\begin{itemize}
\tightlist
\item
  \begin{enumerate}
  \def\labelenumi{(\Alph{enumi})}
  \tightlist
  \item
    AlphaFold needs information from large databases of protein
    sequences and structures\\
  \end{enumerate}
\item
  \begin{enumerate}
  \def\labelenumi{(\Alph{enumi})}
  \setcounter{enumi}{1}
  \tightlist
  \item
    AlphaFold performs \emph{best} on multimeric proteins\\
  \end{enumerate}
\item
  \begin{enumerate}
  \def\labelenumi{(\Alph{enumi})}
  \setcounter{enumi}{2}
  \tightlist
  \item
    AlphaFold is likely to perform poorly on a protein that has a very
    small number of homologs\\
  \end{enumerate}
\item
  \begin{enumerate}
  \def\labelenumi{(\Alph{enumi})}
  \setcounter{enumi}{3}
  \tightlist
  \item
    AlphaFold does not predict the structures of RNA
  \end{enumerate}
\end{itemize}

\textbf{Protein sequence co-evolution gives information on}

\begin{itemize}
\tightlist
\item
  \begin{enumerate}
  \def\labelenumi{(\Alph{enumi})}
  \tightlist
  \item
    The conformations of sidechains in a protein\\
  \end{enumerate}
\item
  \begin{enumerate}
  \def\labelenumi{(\Alph{enumi})}
  \setcounter{enumi}{1}
  \tightlist
  \item
    The presence of tryptophanes in a protein\\
  \end{enumerate}
\item
  \begin{enumerate}
  \def\labelenumi{(\Alph{enumi})}
  \setcounter{enumi}{2}
  \tightlist
  \item
    The stability of a protein\\
  \end{enumerate}
\item
  \begin{enumerate}
  \def\labelenumi{(\Alph{enumi})}
  \setcounter{enumi}{3}
  \tightlist
  \item
    3D contacts between residues in a protein
  \end{enumerate}
\end{itemize}

\textbf{Your target sequence is 90\% identical to a template structure.
You\ldots{} }

\begin{itemize}
\tightlist
\item
  \begin{enumerate}
  \def\labelenumi{(\Alph{enumi})}
  \tightlist
  \item
    Use SwissModel or Modeller, because AlphaFold cannot outperform
    homology modeling with that high sequence similarity\\
  \end{enumerate}
\item
  \begin{enumerate}
  \def\labelenumi{(\Alph{enumi})}
  \setcounter{enumi}{1}
  \tightlist
  \item
    Use Alphafold because the model generation is better\\
  \end{enumerate}
\item
  \begin{enumerate}
  \def\labelenumi{(\Alph{enumi})}
  \setcounter{enumi}{2}
  \tightlist
  \item
    Use Alphafold only if the proteins are orphans\\
  \end{enumerate}
\item
  \begin{enumerate}
  \def\labelenumi{(\Alph{enumi})}
  \setcounter{enumi}{3}
  \tightlist
  \item
    Use I-Tasser
  \end{enumerate}
\end{itemize}

\textbf{Two protein sequences differ by a single amino acid. You use
AlphaFold to predict the structures of those sequences and find two
structure models that are basically identical. You\ldots{} }

\begin{itemize}
\tightlist
\item
  \begin{enumerate}
  \def\labelenumi{(\Alph{enumi})}
  \tightlist
  \item
    Trust those results as this is exactly what you expected\\
  \end{enumerate}
\item
  \begin{enumerate}
  \def\labelenumi{(\Alph{enumi})}
  \setcounter{enumi}{1}
  \tightlist
  \item
    Do not fully trust those results in high detail as single mutations
    are known to often alter the structure of a protein\\
  \end{enumerate}
\item
  \begin{enumerate}
  \def\labelenumi{(\Alph{enumi})}
  \setcounter{enumi}{2}
  \tightlist
  \item
    Trust those results as AlphaFold has solved the protein structure
    prediction problem\\
  \end{enumerate}
\item
  \begin{enumerate}
  \def\labelenumi{(\Alph{enumi})}
  \setcounter{enumi}{3}
  \tightlist
  \item
    Do not trust those results at all as AlphaFold has not solved the
    protein structure prediction problem
  \end{enumerate}
\end{itemize}

\textbf{Which of the following statements is incorrect?}

\begin{itemize}
\tightlist
\item
  \begin{enumerate}
  \def\labelenumi{(\Alph{enumi})}
  \tightlist
  \item
    AlphaFold relies heavily on multiple sequence alignment\\
  \end{enumerate}
\item
  \begin{enumerate}
  \def\labelenumi{(\Alph{enumi})}
  \setcounter{enumi}{1}
  \tightlist
  \item
    AlphaFold is an end-to-end modelling software, i.e., it predicts the
    full conformation of a protein, including the conformations of its
    sidechains\\
  \end{enumerate}
\item
  \begin{enumerate}
  \def\labelenumi{(\Alph{enumi})}
  \setcounter{enumi}{2}
  \tightlist
  \item
    Last version of AlphaFold does not rely on MSAs\\
  \end{enumerate}
\item
  \begin{enumerate}
  \def\labelenumi{(\Alph{enumi})}
  \setcounter{enumi}{3}
  \tightlist
  \item
    AlphaFold does not consider information about ions or ligands when
    reading a file defining the structure of a protein
  \end{enumerate}
\end{itemize}

\textbf{To assess the quality of a protein structure prediction, we
compute either the RMSD or the TM score between the model and the actual
experimental structure of the protein. What can you say about those
scores? }

\begin{itemize}
\tightlist
\item
  \begin{enumerate}
  \def\labelenumi{(\Alph{enumi})}
  \tightlist
  \item
    The higher the RMSD, the better the prediction\\
  \end{enumerate}
\item
  \begin{enumerate}
  \def\labelenumi{(\Alph{enumi})}
  \setcounter{enumi}{1}
  \tightlist
  \item
    The higher the TM score, the better the prediction\\
  \end{enumerate}
\item
  \begin{enumerate}
  \def\labelenumi{(\Alph{enumi})}
  \setcounter{enumi}{2}
  \tightlist
  \item
    A reasonable value for a good TM score is 10
  \end{enumerate}
\end{itemize}

\textbf{To assess the quality of AlphaFold models, it provides the pLDDT
score and the PAE. What can you say about those scores? }

\begin{itemize}
\tightlist
\item
  \begin{enumerate}
  \def\labelenumi{(\Alph{enumi})}
  \tightlist
  \item
    The PAE is used to color the model and it is stored in the PDB
    file\\
  \end{enumerate}
\item
  \begin{enumerate}
  \def\labelenumi{(\Alph{enumi})}
  \setcounter{enumi}{1}
  \tightlist
  \item
    The PAE assess confidence in the domain packing and domain topology
    of the protein\\
  \end{enumerate}
\item
  \begin{enumerate}
  \def\labelenumi{(\Alph{enumi})}
  \setcounter{enumi}{2}
  \tightlist
  \item
    Regions with low pLDDT are expected to be modelled to high accuracy
  \end{enumerate}
\end{itemize}

\textbf{Which of the following best describes the key advancement of
AlphaFold2 compared to the original AlphaFold? }

\begin{itemize}
\tightlist
\item
  \begin{enumerate}
  \def\labelenumi{(\Alph{enumi})}
  \tightlist
  \item
    Introduction of diffusion models\\
  \end{enumerate}
\item
  \begin{enumerate}
  \def\labelenumi{(\Alph{enumi})}
  \setcounter{enumi}{1}
  \tightlist
  \item
    Significantly improved accuracy due to attention mechanisms and
    iterative refinement\\
  \end{enumerate}
\item
  \begin{enumerate}
  \def\labelenumi{(\Alph{enumi})}
  \setcounter{enumi}{2}
  \tightlist
  \item
    Reliance on purely ab initio folding\\
  \end{enumerate}
\item
  \begin{enumerate}
  \def\labelenumi{(\Alph{enumi})}
  \setcounter{enumi}{3}
  \tightlist
  \item
    Ability to predict protein-ligand interactions
  \end{enumerate}
\end{itemize}

\textbf{The ``Invariant Point Attention'' (IPA) mechanism in AlphaFold2
is essential for }

\begin{itemize}
\tightlist
\item
  \begin{enumerate}
  \def\labelenumi{(\Alph{enumi})}
  \tightlist
  \item
    Building a 3D protein structure that is rotationally and
    translationally invariant\\
  \end{enumerate}
\item
  \begin{enumerate}
  \def\labelenumi{(\Alph{enumi})}
  \setcounter{enumi}{1}
  \tightlist
  \item
    Generating multiple sequence alignments\\
  \end{enumerate}
\item
  \begin{enumerate}
  \def\labelenumi{(\Alph{enumi})}
  \setcounter{enumi}{2}
  \tightlist
  \item
    Capturing long-range interactions in the protein sequence\\
  \end{enumerate}
\item
  \begin{enumerate}
  \def\labelenumi{(\Alph{enumi})}
  \setcounter{enumi}{3}
  \tightlist
  \item
    Performing energy minimization
  \end{enumerate}
\end{itemize}

\textbf{How do protein natural language-based models (like ESMFold)
represent proteins? }

\begin{itemize}
\tightlist
\item
  \begin{enumerate}
  \def\labelenumi{(\Alph{enumi})}
  \tightlist
  \item
    As a 3D coordinate map\\
  \end{enumerate}
\item
  \begin{enumerate}
  \def\labelenumi{(\Alph{enumi})}
  \setcounter{enumi}{1}
  \tightlist
  \item
    As a set of torsion angles\\
  \end{enumerate}
\item
  \begin{enumerate}
  \def\labelenumi{(\Alph{enumi})}
  \setcounter{enumi}{2}
  \tightlist
  \item
    As a sequence of amino acid tokens, similar to words in a sentence\\
  \end{enumerate}
\item
  \begin{enumerate}
  \def\labelenumi{(\Alph{enumi})}
  \setcounter{enumi}{3}
  \tightlist
  \item
    As a graph of interacting residues
  \end{enumerate}
\end{itemize}

\chapter{Useful links}\label{links}

\begin{itemize}
\item
  Introductory article to Neural Networks at the IBM site:
  \url{https://www.ibm.com/cloud/learn/neural-networks}
\item
  ColabFold Tutorial presented presented by Sergey Ovchinnikov and
  Martin Steineggerat the Boston Protein Design and Modeling Club (6 ago
  2021). \href{https://www.youtube.com/watch?v=Rfw7thgGTwI}{{[}video{]}}
  \href{https://docs.google.com/presentation/d/1mnffk23ev2QMDzGZ5w1skXEadTe54l8-Uei6ACce8eI}{{[}slides{]}}.
\item
  Post about Alphafold2 in the Oxford Protein Informatics Group site:
  \url{https://www.blopig.com/blog/2021/07/alphafold-2-is-here-whats-behind-the-structure-prediction-miracle/}
\item
  A very good digest article about the AlphaFold2 paper:
  \url{https://moalquraishi.wordpress.com/2021/07/25/the-alphafold2-method-paper-a-fount-of-good-ideas/}
\item
  Post on the AlphaFold revolution meaning in biomedicine at the the UK
  Institute for Cancer Research website:
  \url{https://www.icr.ac.uk/blogs/the-drug-discoverer/page-details/reflecting-on-deepmind-s-alphafold-artificial-intelligence-success-what-s-the-real-significance-for-protein-folding-research-and-drug-discovery}
\item
  A post that explain how AlphaFold and related methods can be used to
  create new structures:
  \url{https://theconversation.com/when-researchers-dont-have-the-proteins-they-need-they-can-get-ai-to-hallucinate-new-structures-173209}
\item
  Selected interesting posts about Alphafold3 modeling process and its
  potential:

  \begin{itemize}
  \item
    \url{https://elanapearl.github.io/blog/2024/the-illustrated-alphafold/}
  \item
    \url{https://carlos.outeiral.net/alphafold-3/}
  \item
    \href{https://research.dimensioncap.com/p/an-opinionated-alphafold3-field-guide}{\textbf{https://research.dimensioncap.com/p/an-opinionated-alphafold3-field-guide}}
  \item
    \url{https://www.labiotech.eu/in-depth/alpha-fold-3-drug-discovery/}
  \end{itemize}
\end{itemize}

\phantomsection\label{refs}
\begin{CSLReferences}{1}{0}
\bibitem[\citeproctext]{ref-abramson2024}
Abramson, Josh, Jonas Adler, Jack Dunger, Richard Evans, Tim Green,
Alexander Pritzel, Olaf Ronneberger, et al. 2024a. {``Accurate structure
prediction of biomolecular interactions with AlphaFold 3.''}
\emph{Nature} 630 (8016): 493--500.
\url{https://doi.org/10.1038/s41586-024-07487-w}.

\bibitem[\citeproctext]{ref-abramson2024a}
---------, et al. 2024b. {``Addendum: Accurate structure prediction of
biomolecular interactions with AlphaFold 3.''} \emph{Nature} 636 (8042):
E4. \url{https://doi.org/10.1038/s41586-024-08416-7}.

\bibitem[\citeproctext]{ref-adiyaman2023}
Adiyaman, Recep, Nicholas S Edmunds, Ahmet G Genc, Shuaa M A Alharbi,
and Liam J McGuffin. 2023. {``Improvement of Protein Tertiary and
Quaternary Structure Predictions Using the ReFOLD Refinement Method and
the AlphaFold2 Recycling Process.''} \emph{Bioinformatics Advances} 3
(1): vbad078. \url{https://doi.org/10.1093/bioadv/vbad078}.

\bibitem[\citeproctext]{ref-aljanabi2022}
Al-Janabi, Aisha. 2022. {``Has DeepMind's AlphaFold Solved the Protein
Folding Problem?''} \emph{BioTechniques} 72 (3): 73--76.
\url{https://doi.org/10.2144/btn-2022-0007}.

\bibitem[\citeproctext]{ref-alphafol2024}
{``AlphaFold3 {\textemdash} Why Did Nature Publish It Without Its
Code?''} 2024. \emph{Nature} 629 (8013): 728--28.
\url{https://doi.org/10.1038/d41586-024-01463-0}.

\bibitem[\citeproctext]{ref-baek}
Baek, Minkyung, Ivan Anishchenko, Ian R. Humphreys, Qian Cong, David
Baker, and Frank DiMaio. n.d. {``Efficient and Accurate Prediction of
Protein Structure Using RoseTTAFold2.''}
\url{https://doi.org/10.1101/2023.05.24.542179}.

\bibitem[\citeproctext]{ref-baek2021}
Baek, Minkyung, Frank DiMaio, Ivan Anishchenko, Justas Dauparas, Sergey
Ovchinnikov, Gyu Rie Lee, Jue Wang, et al. 2021. {``Accurate prediction
of protein structures and interactions using a three-track neural
network.''} \emph{Science (New York, N.Y.)} 373 (6557): 871--76.
\url{https://doi.org/10.1126/science.abj8754}.

\bibitem[\citeproctext]{ref-baker2025}
Baker, David, and Fay Lin. 2025. {``The Promise of Protein Design: A
q\&a with Nobel Laureate David Baker.''} \emph{GEN Biotechnology} 4 (1):
16--18. \url{https://doi.org/10.1089/genbio.2025.0004}.

\bibitem[\citeproctext]{ref-bhattacharya2021}
Bhattacharya, Sutanu, Rahmatullah Roche, Md Hossain Shuvo, and Debswapna
Bhattacharya. 2021. {``Recent Advances in Protein Homology Detection
Propelled by Inter-Residue Interaction Map Threading.''} \emph{Frontiers
in Molecular Biosciences} 8: 643752.
\url{https://doi.org/10.3389/fmolb.2021.643752}.

\bibitem[\citeproctext]{ref-callaway2022}
Callaway, Ewen. 2022. {``AlphaFold{'}s New Rival? Meta AI Predicts Shape
of 600 Million Proteins.''} \emph{Nature} 611 (7935): 211--12.
\url{https://doi.org/10.1038/d41586-022-03539-1}.

\bibitem[\citeproctext]{ref-callaway2024}
---------. 2024. {``Who Will Make AlphaFold3 Open Source? Scientists
Race to Crack AI Model.''} \emph{Nature} 630 (8015): 14--15.
\url{https://doi.org/10.1038/d41586-024-01555-x}.

\bibitem[\citeproctext]{ref-discovery}
Discovery, Chai, Jacques Boitreaud, Jack Dent, Matthew McPartlon, Joshua
Meier, Vinicius Reis, Alex Rogozhnikov, and Kevin Wu. n.d. {``Chai-1:
Decoding the Molecular Interactions of Life.''}
\url{https://doi.org/10.1101/2024.10.10.615955}.

\bibitem[\citeproctext]{ref-elofsson2023}
Elofsson, Arne. 2023. {``Progress at protein structure prediction, as
seen in CASP15.''} \emph{Current Opinion in Structural Biology} 80
(June): 102594. \url{https://doi.org/10.1016/j.sbi.2023.102594}.

\bibitem[\citeproctext]{ref-evans2022}
Evans, Richard, Michael O'Neill, Alexander Pritzel, Natasha Antropova,
Andrew Senior, Tim Green, Augustin Žídek, et al. 2022. {``Protein
Complex Prediction with AlphaFold-Multimer.''}
\url{https://www.biorxiv.org/content/10.1101/2021.10.04.463034v2}.

\bibitem[\citeproctext]{ref-ford}
Ford, Colby T., Samee Ullah, Dinler Amaral Antunes, and Tarsis Gesteira
Ferreira. n.d. {``PyMOLfold: Interactive Protein and Ligand Structure
Prediction in PyMOL.''} \url{https://doi.org/10.48550/arXiv.2502.00508}.

\bibitem[\citeproctext]{ref-garg2025}
Garg, Sriram G., and Georg K. A. Hochberg. 2025. {``A General
Substitution Matrix for Structural Phylogenetics.''} \emph{Molecular
Biology and Evolution} 42 (6): msaf124.
\url{https://doi.org/10.1093/molbev/msaf124}.

\bibitem[\citeproctext]{ref-guo2022a}
Guo, Hao-Bo, Alexander Perminov, Selemon Bekele, Gary Kedziora, Sanaz
Farajollahi, Vanessa Varaljay, Kevin Hinkle, et al. 2022. {``AlphaFold2
models indicate that protein sequence determines both structure and
dynamics.''} \emph{Scientific Reports} 12 (1): 10696.
\url{https://doi.org/10.1038/s41598-022-14382-9}.

\bibitem[\citeproctext]{ref-holm2022}
Holm, Liisa. 2022. {``Dali Server: Structural Unification of Protein
Families.''} \emph{Nucleic Acids Research} 50 (W1): W210--15.
\url{https://doi.org/10.1093/nar/gkac387}.

\bibitem[\citeproctext]{ref-jumper2021}
Jumper, John, Richard Evans, Alexander Pritzel, Tim Green, Michael
Figurnov, Olaf Ronneberger, Kathryn Tunyasuvunakool, et al. 2021.
{``Highly Accurate Protein Structure Prediction with AlphaFold.''}
\emph{Nature} 596 (7873): 583--89.
\url{https://doi.org/10.1038/s41586-021-03819-2}.

\bibitem[\citeproctext]{ref-kempen}
Kempen, Michel van, Stephanie S. Kim, Charlotte Tumescheit, Milot
Mirdita, Cameron L. M. Gilchrist, Johannes Söding, and Martin
Steinegger. n.d. {``Foldseek: Fast and Accurate Protein Structure
Search.''} \url{https://doi.org/10.1101/2022.02.07.479398}.

\bibitem[\citeproctext]{ref-vankempen2023a}
Kempen, Michel van, Stephanie S. Kim, Charlotte Tumescheit, Milot
Mirdita, Jeongjae Lee, Cameron L. M. Gilchrist, Johannes Söding, and
Martin Steinegger. 2023. {``Fast and Accurate Protein Structure Search
with Foldseek.''} \emph{Nature Biotechnology}, May, 1--4.
\url{https://doi.org/10.1038/s41587-023-01773-0}.

\bibitem[\citeproctext]{ref-kumar2025}
Kumar, Ashwani, Aanchal Gupta, Shubham Kumar, Apoorva Haldkar, and Hansa
Wali. 2025. {``Artificial Intelligence (AI)-Based Protein Structure
Prediction and Analysis.''} In, edited by Sudip Mandal, 39--57. New
York, NY: Springer US.
\url{https://doi.org/10.1007/978-1-0716-4690-8_3}.

\bibitem[\citeproctext]{ref-lee}
Lee, Sewon, Gyuri Kim, Eli Levy Karin, Milot Mirdita, Sukhwan Park,
Rayan Chikhi, Artem Babaian, Andriy Kryshtafovych, and Martin
Steinegger. n.d. {``Petascale Homology Search for Structure
Prediction.''} \url{https://doi.org/10.1101/2023.07.10.548308}.

\bibitem[\citeproctext]{ref-lin2023}
Lin, Zeming, Halil Akin, Roshan Rao, Brian Hie, Zhongkai Zhu, Wenting
Lu, Nikita Smetanin, et al. 2023. {``Evolutionary-Scale Prediction of
Atomic-Level Protein Structure with a Language Model.''} \emph{Science}
379 (6637): 1123--30. \url{https://doi.org/10.1126/science.ade2574}.

\bibitem[\citeproctext]{ref-mirdita2022}
Mirdita, Milot, Konstantin Schütze, Yoshitaka Moriwaki, Lim Heo, Sergey
Ovchinnikov, and Martin Steinegger. 2022. {``ColabFold: making protein
folding accessible to all.''} \emph{Nature Methods} 19 (6): 679--82.
\url{https://doi.org/10.1038/s41592-022-01488-1}.

\bibitem[\citeproctext]{ref-moi2025}
Moi, David, Charles Bernard, Martin Steinegger, Yannis Nevers, Mauricio
Langleib, and Christophe Dessimoz. 2025. {``Structural Phylogenetics
Unravels the Evolutionary Diversification of Communication Systems in
Gram-Positive Bacteria and Their Viruses.''} \emph{Nature Structural \&
Molecular Biology} 32 (12): 2492--502.
\url{https://doi.org/10.1038/s41594-025-01649-8}.

\bibitem[\citeproctext]{ref-molodenskiy}
Molodenskiy, Dmitry, Valentin J. Maurer, Dingquan Yu, Grzegorz
Chojnowski, Stefan Bienert, Gerardo Tauriello, Konstantin Gilep, Torsten
Schwede, and Jan Kosinski. n.d. {``AlphaPulldown2 {\textendash} A
General Pipeline for High-Throughput Structural Modeling.''}
\url{https://doi.org/10.1101/2024.11.28.625873}.

\bibitem[\citeproctext]{ref-notin2024}
Notin, Pascal, Nathan Rollins, Yarin Gal, Chris Sander, and Debora
Marks. 2024. {``Machine Learning for Functional Protein Design.''}
\emph{Nature Biotechnology} 42 (2): 216--28.
\url{https://doi.org/10.1038/s41587-024-02127-0}.

\bibitem[\citeproctext]{ref-outeiral}
Outeiral, Carlos, Daniel A. Nissley, and Charlotte M. Deane. n.d.
{``Current Protein Structure Predictors Do Not Produce Meaningful
Folding Pathways.''} \url{https://doi.org/10.1101/2021.09.20.461137}.

\bibitem[\citeproctext]{ref-pakhrin2021}
Pakhrin, Subash C., Bikash Shrestha, Badri Adhikari, and Dukka B. Kc.
2021. {``Deep Learning-Based Advances in Protein Structure
Prediction.''} \emph{International Journal of Molecular Sciences} 22
(11): 5553. \url{https://doi.org/10.3390/ijms22115553}.

\bibitem[\citeproctext]{ref-pederson2021}
Pederson, Thoru. 2021. {``Protein Structure: Has Levinthal's Paradox
{``}Folded{''}?''} \emph{The FASEB Journal} 35 (3): e21416.
\url{https://doi.org/10.1096/fj.202100136}.

\bibitem[\citeproctext]{ref-peng2023}
Peng, Zhenling, Wenkai Wang, Hong Wei, Xiaoge Li, and Jianyi Yang. 2023.
{``Improved protein structure prediction with trRosettaX2, AlphaFold2,
and optimized MSAs in CASP15.''} \emph{Proteins}, August.
\url{https://doi.org/10.1002/prot.26570}.

\bibitem[\citeproctext]{ref-portapardo2022}
Porta-Pardo, Eduard, Victoria Ruiz-Serra, Samuel Valentini, and Alfonso
Valencia. 2022. {``The Structural Coverage of the Human Proteome Before
and After AlphaFold.''} \emph{PLOS Computational Biology} 18 (1):
e1009818. \url{https://doi.org/10.1371/journal.pcbi.1009818}.

\bibitem[\citeproctext]{ref-puente-lelievre}
Puente-Lelievre, Caroline, Ashar J. Malik, Jordan Douglas, David Ascher,
Matthew Baker, Jane Allison, Anthony Poole, et al. n.d.
{``Tertiary-Interaction Characters Enable Fast, Model-Based Structural
Phylogenetics Beyond the Twilight Zone.''}
\url{https://doi.org/10.1101/2023.12.12.571181}.

\bibitem[\citeproctext]{ref-rouger}
Rouger, Quentin, Emmanuel Giudice, Damien F. Meyer, and Kévin Macé. n.d.
{``PPIFold: A Tool for Analysis of Protein-Protein Interaction from
AlphaPullDown.''} \url{https://doi.org/10.1101/2025.01.08.631489}.

\bibitem[\citeproctext]{ref-senior2020}
Senior, Andrew W., Richard Evans, John Jumper, James Kirkpatrick,
Laurent Sifre, Tim Green, Chongli Qin, et al. 2020. {``Improved protein
structure prediction using potentials from deep learning.''}
\emph{Nature} 577 (7792): 706--10.
\url{https://doi.org/10.1038/s41586-019-1923-7}.

\bibitem[\citeproctext]{ref-shin2025}
Shin, Heewhan, Alexandria Holland, Abdulrazak Alsaleh, Alyssa D. Retiz,
Ying Z. Pigli, Oluwateniola T. Taiwo-Aiyerin, Tania Peña Reyes, et al.
2025. {``Identification of cognate recombination directionality factors
for large serine recombinases by virtual pulldown.''} \emph{Nucleic
Acids Research} 53 (14): gkaf691.
\url{https://doi.org/10.1093/nar/gkaf691}.

\bibitem[\citeproctext]{ref-team}
Team, ByteDance AML AI4Science, Xinshi Chen, Yuxuan Zhang, Chan Lu,
Wenzhi Ma, Jiaqi Guan, Chengyue Gong, et al. n.d. {``Protenix -
Advancing Structure Prediction Through a Comprehensive AlphaFold3
Reproduction.''} \url{https://doi.org/10.1101/2025.01.08.631967}.

\bibitem[\citeproctext]{ref-varadi2022}
Varadi, Mihaly, Stephen Anyango, Mandar Deshpande, Sreenath Nair, Cindy
Natassia, Galabina Yordanova, David Yuan, et al. 2022. {``AlphaFold
Protein Structure Database: Massively Expanding the Structural Coverage
of Protein-Sequence Space with High-Accuracy Models.''} \emph{Nucleic
Acids Research} 50 (D1): D439--44.
\url{https://doi.org/10.1093/nar/gkab1061}.

\bibitem[\citeproctext]{ref-wang2022}
Wang, Wenkai, Zhenling Peng, and Jianyi Yang. 2022. {``Single-sequence
protein structure prediction using supervised transformer protein
language models.''} \emph{Nature Computational Science} 2 (12): 804--14.
\url{https://doi.org/10.1038/s43588-022-00373-3}.

\bibitem[\citeproctext]{ref-wankowicz2024}
Wankowicz, Stephanie, Pedro Beltrao, Benjamin Cravatt, Roland Dunbrack,
Anthony Gitter, Kresten Lindorff-Larsen, Sergey Ovchinnikov, Nicholas
Polizzi, Brian Shoichet, and James Fraser. 2024. {``AlphaFold3
Transparency and Reproducibility,''} May.
\url{https://zenodo.org/records/11391920}.

\bibitem[\citeproctext]{ref-wilson2022}
Wilson, Carter J., Wing-Yiu Choy, and Mikko Karttunen. 2022.
{``AlphaFold2: A Role for Disordered Protein/Region Prediction?''}
\emph{International Journal of Molecular Sciences} 23 (9): 4591.
\url{https://doi.org/10.3390/ijms23094591}.

\bibitem[\citeproctext]{ref-wohlwend}
Wohlwend, Jeremy, Gabriele Corso, Saro Passaro, Mateo Reveiz, Ken
Leidal, Wojtek Swiderski, Tally Portnoi, et al. n.d. {``Boltz-1:
Democratizing Biomolecular Interaction Modeling.''}
\url{https://doi.org/10.1101/2024.11.19.624167}.

\bibitem[\citeproctext]{ref-yang2020}
Yang, Jianyi, Ivan Anishchenko, Hahnbeom Park, Zhenling Peng, Sergey
Ovchinnikov, and David Baker. 2020. {``Improved Protein Structure
Prediction Using Predicted Interresidue Orientations.''}
\emph{Proceedings of the National Academy of Sciences} 117 (3):
1496--1503. \url{https://doi.org/10.1073/pnas.1914677117}.

\bibitem[\citeproctext]{ref-yu2023}
Yu, Dingquan, Grzegorz Chojnowski, Maria Rosenthal, and Jan Kosinski.
2023. {``AlphaPulldown{\textemdash}a Python Package for
Protein{\textendash}protein Interaction Screens Using
AlphaFold-Multimer.''} \emph{Bioinformatics} 39 (1): btac749.
\url{https://doi.org/10.1093/bioinformatics/btac749}.

\bibitem[\citeproctext]{ref-zheng2021}
Zheng, Wei, Chengxin Zhang, Yang Li, Robin Pearce, Eric W. Bell, and
Yang Zhang. 2021. {``Folding non-homologous proteins by coupling
deep-learning contact maps with I-TASSER assembly simulations.''}
\emph{Cell Reports Methods} 1 (3): 100014.
\url{https://doi.org/10.1016/j.crmeth.2021.100014}.

\end{CSLReferences}


\backmatter


\end{document}
